\documentclass[12pt,a4paper,oneside]{report}             % Single-side
%\documentclass[11pt,a4paper,twoside,openright]{report}  % Duplex

%\PassOptionsToPackage{chapternumber=Huordinal}{magyar.ldf}

\usepackage{fontspec}
\setmainfont{Times New Roman}  % specify a font that supports required glyphs

\usepackage{amsmath}
\usepackage{amssymb}
\usepackage[english,magyar]{babel}
\usepackage{hyperref}
\usepackage{datetime}
\usepackage{enumerate}
\usepackage[thmmarks]{ntheorem}
\usepackage{graphics}
\usepackage{epsfig}
\usepackage{listings}
\usepackage{color}
%\usepackage{fancyhdr}
\usepackage{lastpage}
\usepackage{anysize}
\usepackage{sectsty}
\usepackage{setspace}  % Ettol a tablazatok, abrak, labjegyzetek maradnak 1-es sorkozzel!
\usepackage[hang]{caption}
\usepackage{pdfpages}
\usepackage{todonotes}
\usepackage{enumitem}
\usepackage{makecell}


%--------------------------------------------------------------------------------------
% Main variables
%--------------------------------------------------------------------------------------
\newcommand{\vikszerzo}{Nyíri Tamás}
\newcommand{\vikkonzulens}{Dr.~Szeberényi Imre}
\newcommand{\vikcim}{Feladatellenőrző rendszer továbbfejlesztése}
\newcommand{\viktanszek}{Irányítástechnika és Informatika Tanszék}
\newcommand{\vikdoktipus}{Szakdolgozat}

%--------------------------------------------------------------------------------------
% Page layout setup
%--------------------------------------------------------------------------------------
% we need to redefine the pagestyle plain
% another possibility is to use the body of this command without \fancypagestyle
% and use \pagestyle{fancy} but in that case the special pages
% (like the ToC, the References, and the Chapter pages)remain in plane style

\pagestyle{plain}
%\setlength{\parindent}{0pt} % áttekinthetőbb, angol nyelvű dokumentumokban jellemző
%\setlength{\parskip}{8pt plus 3pt minus 3pt} % áttekinthetőbb, angol nyelvű dokumentumokban jellemző
\setlength{\parindent}{12pt} % magyar nyelvű dokumentumokban jellemző
\setlength{\parskip}{0pt}    % magyar nyelvű dokumentumokban jellemző

\marginsize{35mm}{25mm}{15mm}{15mm} % anysize package
\setcounter{secnumdepth}{0}
\sectionfont{\large\upshape\bfseries}
\setcounter{secnumdepth}{2}
\singlespacing
\frenchspacing

%--------------------------------------------------------------------------------------
%	Setup hyperref package
%--------------------------------------------------------------------------------------
\hypersetup{
    bookmarks=true,            % show bookmarks bar?
    unicode=true,              % non-Latin characters in Acrobat’s bookmarks
    pdftitle={\vikcim},        % title
    pdfauthor={\vikszerzo},    % author
    pdfsubject={\vikdoktipus}, % subject of the document
    pdfcreator={\vikszerzo},   % creator of the document
    pdfproducer={Producer},    % producer of the document
    pdfkeywords={keywords},    % list of keywords
    pdfnewwindow=true,         % links in new window
    colorlinks=true,           % false: boxed links; true: colored links
    linkcolor=black,           % color of internal links
    citecolor=black,           % color of links to bibliography
    filecolor=black,           % color of file links
    urlcolor=black             % color of external links
}

%--------------------------------------------------------------------------------------
% Set up listings
%--------------------------------------------------------------------------------------
\lstset{
	basicstyle=\scriptsize\ttfamily, % print whole listing small
	keywordstyle=\color{black}\bfseries\underbar, % underlined bold black keywords
	identifierstyle=, 					% nothing happens
	commentstyle=\color{white}, % white comments
	stringstyle=\scriptsize\sffamily, 			% typewriter type for strings
	showstringspaces=false,     % no special string spaces
	aboveskip=3pt,
	belowskip=3pt,
	columns=fixed,
	backgroundcolor=\color{lightgray},
} 		
\def\lstlistingname{lista}	

%--------------------------------------------------------------------------------------
%	Some new commands and declarations
%--------------------------------------------------------------------------------------
\newcommand{\code}[1]{{\upshape\ttfamily\scriptsize\indent #1}}

% define references
\newcommand{\figref}[1]{\ref{figure:#1}. ábra}
\renewcommand{\eqref}[1]{(\ref{eq:#1})}
\newcommand{\listref}[1]{\ref{listing:#1}.}
\newcommand{\sectref}[1]{\ref{sect:#1}}
\newcommand{\tabref}[1]{\ref{tab:#1}.}

\newcommand{\sectsep}{
	\par \nopagebreak
	\parbox{\linewidth}{
    	\centering\bigskip$\ast\quad\ast\quad\ast$\par\bigskip
	}
}

\DeclareMathOperator*{\argmax}{arg\,max}
%\DeclareMathOperator*[1]{\floor}{arg\,max}
\DeclareMathOperator{\sign}{sgn}
\DeclareMathOperator{\rot}{rot}
\definecolor{lightgray}{rgb}{0.95,0.95,0.95}

\author{\vikszerzo}
\title{\viktitle}
\includeonly{
	abstract,
	acknowledgement,
	bibliography,
	declaration,
	summary,
	titlepage,
}
%--------------------------------------------------------------------------------------
%	Setup captions
%--------------------------------------------------------------------------------------
\captionsetup[figure]{
%labelsep=none,
%font={footnotesize,it},
%justification=justified,
width=.75\textwidth,
aboveskip=10pt}

\renewcommand{\captionlabelfont}{\small\bf}
\renewcommand{\captionfont}{\footnotesize\it}

%--------------------------------------------------------------------------------------
% Table of contents and the main text
%--------------------------------------------------------------------------------------
\begin{document}
\pagenumbering{arabic}
\spacing{1.5} %\onehalfspacing

%--------------------------------------------------------------------------------------
%	The title page
%--------------------------------------------------------------------------------------
\begin{titlepage}
\begin{center}
\includegraphics[width=60mm,keepaspectratio]{figures/BMElogo.png}\\
\vspace{0.3cm}
\textbf{Budapesti Műszaki és Gazdaságtudományi Egyetem}\\
\textmd{Villamosmérnöki és Informatikai Kar}\\
\textmd{\viktanszek}\\[5cm]

\vspace{0.4cm}
{\huge \bfseries \vikcim}\\[0.8cm]
\vspace{0.5cm}
\textsc{\Large \vikdoktipus}\\[4cm]

\begin{tabular}{cc}
 \makebox[7cm]{\emph{Készítette}} & \makebox[7cm]{\emph{Konzulens}} \\
 \makebox[7cm]{\vikszerzo} & \makebox[7cm]{\vikkonzulens}
\end{tabular}

\vfill
{\large \today}
\end{center}
\end{titlepage}



\includepdf[pages={1}]{feladatkiiras.pdf}
%--------------------------------------------------------------------------------------
% Nyilatkozat
%--------------------------------------------------------------------------------------
\begin{center}
\large
\textbf{HALLGATÓI NYILATKOZAT}\\
\end{center}

Alulírott \emph{\vikszerzo}, szigorló hallgató kijelentem, hogy ezt a szakdolgozatot meg nem engedett segítség nélkül, saját magam készítettem, csak a megadott forrásokat (szakirodalom, eszközök stb.) használtam fel. Minden olyan részt, melyet szó szerint, vagy azonos értelemben, de átfogalmazva más forrásból átvettem, egyértelműen, a forrás megadásával megjelöltem.

Hozzájárulok, hogy a jelen munkám alapadatait (szerző, cím, angol és magyar nyelvű tartalmi kivonat, készítés éve, konzulens neve) a BME VIK nyilvánosan hozzáférhető elektronikus formában, a munka teljes szövegét pedig az egyetem belső hálózatán keresztül (vagy autentikált felhasználók számára) közzétegye. Kijelentem, hogy a benyújtott munka és annak elektronikus verziója megegyezik. Dékáni engedéllyel titkosított diplomatervek esetén a dolgozat szövege csak 3 év eltelte után válik hozzáférhetővé.

\begin{flushleft}
\vspace*{1cm}
Budapest, \today
\end{flushleft}

\begin{flushright}
 \vspace*{1cm}
 \makebox[7cm]{\rule{6cm}{.4pt}}\\
 \makebox[7cm]{\emph{\vikszerzo}}\\
 \makebox[7cm]{hallgató}
\end{flushright}
\thispagestyle{empty}

\vfill
\clearpage
\thispagestyle{empty} % an empty page


\tableofcontents\vfill
%----------------------------------------------------------------------------
% Abstract in hungarian
%----------------------------------------------------------------------------
\chapter*{Kivonat}%\addcontentsline{toc}{chapter}{Kivonat}


\vfill

%----------------------------------------------------------------------------
% Abstract in english
%----------------------------------------------------------------------------
\begin{otherlanguage}{english}
\chapter*{Abstract}%\addcontentsline{toc}{chapter}{Abstract}


\end{otherlanguage}
\vfill


%----------------------------------------------------------------------------
\chapter*{Bevezető}\addcontentsline{toc}{chapter}{Bevezető}
%----------------------------------------------------------------------------


%----------------------------------------------------------------------------
\chapter*{Összefoglalás}\addcontentsline{toc}{chapter}{Összefoglalás}
%----------------------------------------------------------------------------

A szakdolgozatom keretében részletesen megismertem mind a JPorta rendszerét, mind a felé támasztott elvárásokat oktatói és hallgatói szempontból. Ezek alapján főként az aktuálisan legkritikusabbnak ítélt területek fejlesztésével foglalkoztam, azaz a dinamikus mezőkkel és az automatikusan kiértékelődő feladatok rendszerével.

Előbbin elért eredményeknek köszönhetően az oktatók már létrehozhatnak olyan dinamikus értékelési mezőket, melyek más dinamikus mezőket is használnak a számításaik során. Ezzel segítik az oktatókat a bonyolultabb mezők elkészítésében. Tipikusan ez a félév végi eredmény meghatározásánál lehet különösen hasznos.

Az automatikus kiértékelő rendszer fejlesztésének köszönhetően már lehetőség van oktatóknak és hallgatóknak is jogosultságot adni a feladatok létrehozására. Ez egy további lépést jelent a tárgyak széleskörű támogatására, hiszen így a feladatok elkészítésére nem csak a portál adminisztrátorait lehet megkérni, hanem minden arra jogosult személy kísérletezhet.

A kódlefedettség ellenőrző pedig lehetőséget teremt a C és C++ nyelvű feladatok esetén a tesztesetek értékelésére. Így könnyebb automatikusan kiszűrni a hibákat, működési zavarokat a beadott megoldásokban.

Végezetül pedig további hasznos funkciókra tettem javaslatot, melyek tovább növelhetik az oktatók és hallgatók elégedettségét a portállal szemben.
%----------------------------------------------------------------------------
\chapter*{Köszönetnyilvánítás}\addcontentsline{toc}{chapter}{Köszönetnyilvánítás}
%----------------------------------------------------------------------------

Szerenték köszönetet mondani konzulensemnek, Dr.~Szeberényi Imrének, aki a tanulmányaim és szakdolgozatom készítésének során végig támogatott.

Emellett külön köszönet jár a BME Informatikai Központ munkatársainak és volt munkatársainak, köztük Dudás Ádámnak és Kálmán Viktornak, akik a korábbi években elkészítették a JPorta fő moduljait és ezekkel kapcsolatban a félév során hasznos információkkal segítették munkámat.
\listoffigures\addcontentsline{toc}{chapter}{Ábrák jegyzéke}
%\listoftables\addcontentsline{toc}{chapter}{Táblázatok jegyzéke}
%\printglossary[title=Rövidítések jegyzéke]

\bibliography{bibliography}\addcontentsline{toc}{chapter}{Irodalomjegyzék}
\bibliographystyle{huplain}

\label{page:last}
\end{document}
